\chapter{Introdução}\label{cap:introducao}
O estudo da Matemática no Ensino Básico brasileiro, essencial para o desenvolvimento do pensamento lógico e da capacidade de análise crítica, enfrenta muitos desafios, dentre eles, um dos mais persistentes: o desinteresse e a descontextualização da disciplina. Para muitos estudantes, a Matemática é percebida como um conjunto de fórmulas abstratas, o que contribui para o baixo desempenho e para a visão de que a matéria é ``chata'' e pouco relevante para a vida prática. Essa crítica à abordagem tradicional é defendida por \citeonline{dambrosio1999etnomatematica}, que argumenta que a ênfase excessiva em um modelo eurocêntrico e a-histórico resulta em uma ``matemática obsoleta e desinteressante'' para o aluno, tornando imprescindível que o ensino da matemática busque metodologias que resgatem o significado e a atratividade do saber.

A Geometria, um campo que exige a visualização espacial e a compreensão rigorosa de relações complexas, é particularmente vulnerável a essa crise. O abandono do ensino da Geometria no Brasil é um fato histórico e curricular amplamente documentado por \citeonline{pavanello1993abandono}. Segundo a autora, a falha reside na priorização de um ensino excessivamente formal e abstrato, onde a didática se concentra na memorização de propriedades e na apresentação dos teoremas exclusivamente de forma escrita e robusta. Esse modelo, ao não fornecer o suporte visual necessário, impede a transição do pensamento intuitivo para o pensamento dedutivo eficaz, transformando o rigor das demonstrações geométricas em um obstáculo intransponível. \citeonline{pavanello1993abandono}, também conclui que o trabalho com a Geometria, quando bem conduzido, deve favorecer o desenvolvimento de um pensamento crítico e autônomo. É nesse contexto de desengajamento e falha na transposição didática que a inovação metodológica se torna uma solução viável para resgatar o valor da disciplina.

%% ------- JUSTIFICATIVCA PARA O TEMA -- INICIO  

\section{Justificativa para a escolha do tema}

A escolha do tema desta dissertação, o ensino de demonstrações em Geometria Plana por meio de um material que combina texto rigoroso e vídeos curtos de animação, justifica-se pela necessidade de intervir diretamente na crise de interesse e nas dificuldades de visualização e de rigor apontadas por \citeonline{pavanello1993abandono}. Em vez de propor mais um material estático, este trabalho organiza o recurso em duas camadas complementares: um vídeo curto, de formato breve e atrativo, que desperta o interesse e dá acesso ao tema, e um texto rigoroso, que conduz o leitor até a demonstração formal. A primeira camada aproxima; a segunda aprofunda. Para fundamentar esse desenho, o trabalho organiza-se em torno de quatro frentes, cada uma respondendo a uma necessidade do problema exposto: a condução do raciocínio próprio da demonstração, a progressão ordenada até o rigor dedutivo, a combinação eficaz de texto e imagem e a aproximação do tema a um público amplo. Essas frentes correspondem, respectivamente, à filosofia da resolução de problemas, à progressão do pensamento geométrico, ao desenho da instrução multimídia e à divulgação matemática, detalhadas a seguir. 

A primeira frente trata a demonstração não como um resultado a memorizar, mas como um processo de descoberta a compreender, priorizando o desenvolvimento do raciocínio sobre a memorização. Nesse sentido, a resolução de problemas é o veículo para o desenvolvimento da heurística, essencial para que o estudante se torne um resolvedor autônomo. Essa postura é endossada por \citeonline{polya2006arte}, que enfatiza a importância da prática e da descoberta pessoal:

\begin{citacao}
Uma grande descoberta resolve um grande problema, mas há sempre uma pitada de descoberta na resolução de qualquer problema. Se ele desafiar a curiosidade e puser em jogo as faculdades inventivas, quem o resolver por seus próprios meios, experimentará a tensão e gozará o triunfo da descoberta \cite[p.~V]{polya2006arte}.
\end{citacao}

O material, portanto, apoia-se nas etapas do método de resolução de problemas, conduzindo o leitor da compreensão do problema à elaboração e à execução de um plano de solução. Mais do que verificar, a demonstração deve explicar, e é essa compreensão do raciocínio que o recurso busca promover.

A segunda frente é a progressão do pensamento geométrico. O rigor da Geometria exige que o estudante avance por etapas no seu entendimento das formas. Este trabalho adota a teoria de \citeonline{vanhiele1986structure}, que descreve o desenvolvimento do pensamento geométrico em níveis hierárquicos, para orientar a organização dos conteúdos. As animações atuam como o suporte visual que apoia a transição do reconhecimento intuitivo para a dedução formal, justamente o ponto em que o material puramente escrito costuma falhar.

A terceira frente é o desenho da instrução multimídia. A combinação de texto e animação favorece a aprendizagem, conforme a Teoria Cognitiva da Aprendizagem Multimídia de \citeonline{mayer2021multimedia}, que trata do gerenciamento da capacidade limitada da memória de trabalho, com frequência sobrecarregada por demonstrações longas e puramente textuais. É o que fundamenta a exibição do texto ao lado do vídeo e a divisão do conteúdo em passos curtos.

A quarta frente é a divulgação matemática. Para além da aprendizagem, o vídeo curto, no estilo que circula nas redes sociais, cumpre um papel comunicativo: seu formato breve e de baixa barreira de entrada aproxima a Geometria de um público amplo. O vídeo atrai e o texto aprofunda, e é essa combinação de funções que distingue a proposta de um material puramente escrito.


\section{Revisão Bibliográfica}
%%%% --------------------
O primeiro pilar desta revisão é o desenvolvimento do pensamento geométrico. O modelo de \citeonline{vanhiele1986structure} é a referência principal para o desenho didático do material, pois diagnostica uma falha central na transposição da Geometria no Ensino Básico: a incompatibilidade entre o nível de raciocínio do aluno e a complexidade do material apresentado. A teoria postula que o pensamento evolui através de níveis hierárquicos, sendo difícil ao estudante ascender a um nível superior sem a assimilação do anterior. O rigor das demonstrações geométricas corresponde ao nível de dedução formal, e a dificuldade em alcançá-lo é apontada pela literatura nacional. Conforme observaram \citeonline[p.~3]{kaleff1994desenvolvimento}, ao investigarem o modelo no contexto brasileiro: ``Seus trabalhos revelam uma alarmante falta de harmonia entre o ensino e aprendizado em matemática''.

Para este trabalho, a teoria de \citeonline{vanhiele1986structure} é utilizada para orientar a sequência dos conteúdos. O material é, portanto, uma intervenção didática intencional que emprega o recurso multimídia como apoio à transição entre os níveis, do reconhecimento intuitivo até a prova rigorosa.

O segundo pilar é a filosofia da resolução de problemas. O material trata a demonstração como veículo para o desenvolvimento da heurística, essencial para que o estudante se torne um resolvedor autônomo. Esse processo apoia-se no método de \citeonline{polya2006arte}, cujas etapas conduzem o leitor da compreensão do problema à elaboração e à execução de um plano de solução, incentivando ativamente a descoberta e a invenção.

O terceiro pilar é o Desenho da Instrução Multimídia. A proposta de utilizar vídeos de 3 minutos, programados em JavaScript e integrados ao texto, é a resposta metodológica à falha do material estático. Esta metodologia é validada pela Teoria Cognitiva da Aprendizagem Multimídia (TCAM) de \citeonline{mayer2021multimedia}. O principal foco da TCAM é o gerenciamento da Capacidade Limitada da memória de trabalho do aluno, que tende a ser sobrecarregada por demonstrações complexas e puramente textuais. Este material visa aplicar diretamente três princípios centrais:

\begin{itemize}
    \item O Princípio Multimídia justifica a combinação do texto com a animação (imagem), pois as pessoas aprendem melhor quando se combina palavras e imagens do que apenas palavras \cite{mayer2021multimedia}.
    \item O Princípio da Contiguidade Temporal sustenta a exibição do texto ao lado do vídeo, garantindo que o aluno processe o que vê e o que ouve simultaneamente.

    \item O Princípio da Segmentação valida o formato de vídeos curtos (3 minutos), pois a divisão do conteúdo em passos menores e controláveis evita a sobrecarga cognitiva (Capacidade Limitada).
\end{itemize}

O quarto pilar é a divulgação matemática. Enquanto a aprendizagem multimídia explica por que, uma vez diante do material, o estudante aprende, a divulgação explica por que o formato atrai e amplia o acesso, incluindo o público no debate sobre o conhecimento \cite{bueno2010comunicacao}. As redes sociais, acessadas por bilhões de pessoas sobretudo por dispositivos móveis, têm reconhecido potencial para a divulgação da ciência \cite{chagasmassarani2020manual}. É nesse ambiente que se inscreve o vídeo curto vertical: de baixa barreira de entrada, ele desperta o interesse de um público amplo e o conduz ao aprofundamento oferecido pelo texto. Essa camada de engajamento, articulada à camada didática, constitui o modelo de duas etapas que caracteriza a proposta.

%%%%---------------------


% ------ OBJETIVOS 

\section{Objetivos}
\subsection{Objetivos Gerais}
O presente trabalho tem como objetivo geral desenvolver e documentar um modelo de material didático-multimídia para o ensino de demonstrações em Geometria Plana, combinando texto rigoroso e vídeos curtos de animação, instanciado em um módulo inicial, e analisar seu potencial teórico-didático à luz da aprendizagem multimídia e da divulgação matemática.
\subsection{Objetivos específicos}

Para a consecução do objetivo geral, esta dissertação propõe-se a:

\begin{enumerate}
    \item Reunir e formalizar, com o rigor adequado, as demonstrações dos principais teoremas e corolários de um módulo inicial de Geometria Plana.

    \item Organizar a progressão didática do material com base nos níveis de pensamento geométrico de \citeonline{vanhiele1986structure}, empregando as animações como suporte visual para a transição ao pensamento dedutivo formal.

    \item Desenvolver e documentar a metodologia de produção dos vídeos curtos de animação, detalhando as etapas de construção, \textit{storyboard} e narração.

    \item Analisar o design do Produto Educacional à luz dos princípios da Teoria Cognitiva da Aprendizagem Multimídia de \citeonline{mayer2021multimedia} e do campo da divulgação matemática, discutindo a complementaridade entre a camada de engajamento (o vídeo) e a camada de aprofundamento (o texto).

    \item Apresentar o Produto Educacional, composto pelo módulo inicial da trilha e seus vídeos, de forma documentada e disponível publicamente, indicando os desdobramentos futuros, como a ampliação para novos módulos.
\end{enumerate}

 %%% ------- PROCEDIMENTOS METODOLOGICOS
 
\section{Procedimentos metodológicos}
A presente dissertação configura-se como uma Pesquisa Bibliográfica e de Desenvolvimento Teórico, cujo foco principal é a criação e a fundamentação de um Produto Educacional (a Trilha Multimídia) para intervir em um problema didático da prática.

\subsection{Abordagem e Natureza}
\begin{itemize}
    \item Natureza da Pesquisa: O trabalho possui a natureza de Pesquisa Bibliográfica de Desenvolvimento Teórico. Ele se alinha à definição de Pesquisa Exploratória de \citeonline{gerhardt2009metodos}, que é realizada quando o objetivo é proporcionar maior familiaridade com um problema, buscando referências para a criação de um modelo teórico e didático.
    \item Abordagem: O trabalho adota uma abordagem qualitativa, uma vez que visa à análise interpretativa das características, do design e das potencialidades do Produto Educacional desenvolvido. Essa escolha metodológica se justifica porque, conforme \citeonline{ludke1986pesquisa}, na pesquisa qualitativa o pesquisador é o principal instrumento de investigação, atuando diretamente na interpretação do objeto em seu contexto, o que é fundamental para alcançar os objetivos propostos.
\end{itemize}

\subsection{Objetivos e Procedimentos}
\begin{itemize}
    \item Objetivos da Pesquisa: O estudo é de natureza Exploratório-Descritiva. É exploratório ao buscar na literatura a crítica ao ensino da Geometria e as soluções multimídia. É descritivo ao detalhar a estrutura, o conteúdo e a metodologia de design da Trilha de Estudos e dos vídeos de animação \cite{gerhardt2009metodos}.
    \item Procedimentos Técnicos: O trabalho baseia-se primordialmente na Pesquisa Bibliográfica. Foram analisadas obras e artigos de referência nos campos da Didática da Matemática, da Tecnologia Educacional e da historiografia do ensino da Geometria. Esta pesquisa embasa todas as decisões tomadas no design do produto.
\end{itemize}

%%% ----- PRODUTO EDUCACIONAL

\section{Produto Educacional}

A culminação desta pesquisa bibliográfica e metodológica é o Produto Educacional, que se materializa na forma de uma trilha de estudos digital em Geometria Plana, publicada como página na internet (o sítio \textit{Geometria Demonstrada}). Nesta etapa, a trilha é apresentada como um módulo inicial, voltado aos conteúdos fundamentais da Geometria Plana, concebido de modo a ser ampliado em novos módulos. Desenvolvido sob as diretrizes do Mestrado Profissional, o recurso visa ser um artefato didático de alta replicabilidade. \\
O Recurso Educacional é composto por:

\begin{enumerate}
    \item Material Escrito (Texto e Estrutura): Consiste nas demonstrações rigorosas dos teoremas e nos exercícios resolvidos, organizados sequencialmente conforme os níveis de pensamento geométrico de \citeonline{vanhiele1986structure}. É a camada de aprofundamento, destinada a quem deseja compreender a demonstração em sua essência lógica.

\item Material Audiovisual: É a camada de engajamento do projeto. Consiste em uma série de vídeos curtos de animação (cerca de três minutos), no estilo \textit{Reels}/\textit{TikTok}, que ilustram as etapas mais complexas das demonstrações geométricas. Pelo formato breve e de baixa barreira de entrada, esses vídeos têm potencial de divulgação, aproximando a Geometria de um público amplo. Conforme a Figura \ref{fig:flame-de-animacao}, que ilustra um quadro da animação, o recurso é desenhado para guiar visualmente o leitor pelas construções auxiliares necessárias à prova dos teoremas.
\begin{figure}[!h]
    \caption{flames de animação: (a) bissetriz interna, (b) ENQ 2025} 
    \vspace{0.25cm}
	\label{fig:flame-de-animacao}
	\centering
	\subfloat[]{		\includegraphics[scale=0.13]{figuras/print1.jpg}
	}\hspace{0.2cm} 
	\subfloat[]{
		\includegraphics[scale=0.13]{figuras/print2.jpg}
	}	
	\fonte{Autoria própria (2025)}
\end{figure}

\end{enumerate}

A produção deste Material Audiovisual é realizada através de um software de animação especializado, que agiliza a criação do conteúdo multimídia. Esta escolha metodológica garante a alta qualidade visual, a rapidez na produção do material e a capacidade de focar o esforço do pesquisador no design didático, e não na codificação bruta.

A trilha é disponibilizada em uma plataforma digital de acesso público, que integra o texto e os vídeos de suporte lado a lado, conforme o princípio da contiguidade de \citeonline{mayer2021multimedia}. É nessa integração que as duas camadas se encontram: o vídeo atrai, o texto aprofunda.



\section{Estrutura da Dissertação}
Este trabalho está organizado em cinco capítulos, além do Produto Educacional, que é a peça central da pesquisa. A divisão é a seguinte:

\begin{itemize}
    \item Capítulo 1 – Introdução: apresenta o contexto da crise de interesse e do abandono do ensino da Geometria, o problema de pesquisa, os objetivos, a relevância e a justificativa para a escolha do tema (demonstração em Geometria, multimídia e divulgação).
    \item Capítulo 2 – Fundamentação Teórica: detalha os pilares conceituais que sustentam a proposta, organizados em dois movimentos complementares. De um lado, a camada de divulgação (o vídeo curto como meio de engajamento); de outro, a camada didática (a progressão do pensamento geométrico de Van Hiele e a resolução de problemas de Pólya). A Teoria Cognitiva da Aprendizagem Multimídia de Mayer articula os dois, explicando por que a combinação de texto e vídeo funciona.
    \item Capítulo 3 – Metodologia: descreve o design da pesquisa, de natureza teórica e de desenvolvimento, detalhando o processo de elaboração do Produto Educacional, a definição do público e os procedimentos de criação das animações.
    \item Capítulo 4 – O Produto Educacional: apresenta, em detalhes, o módulo inicial da trilha de Geometria Plana. Este capítulo contém o cerne do trabalho, incluindo as demonstrações formais, os exercícios resolvidos e a análise teórica do design do material à luz da fundamentação.
    \item Capítulo 5 – Considerações Finais e Trabalhos Futuros: sintetiza os resultados, verifica o cumprimento dos objetivos e aponta desdobramentos, em especial a ampliação do produto para novos módulos.
\end{itemize}

O Produto Educacional, composto pelo módulo inicial da trilha de Geometria Plana e seus vídeos de animação, é a materialização desta pesquisa.
